% SPEC-SW-006 --- Multi-Agent Simulation Engine
% Shared preamble for all spec documents.
% Include from each spec: % Shared preamble for all spec documents.
% Include from each spec: % Shared preamble for all spec documents.
% Include from each spec: \input{../preamble.tex} (adjust depth).

\documentclass[11pt,a4paper]{article}

\usepackage[utf8]{inputenc}
\usepackage[T1]{fontenc}
\usepackage{lmodern}
\usepackage[a4paper,margin=2.3cm]{geometry}
\usepackage{parskip}
\usepackage{microtype}
\usepackage[dvipsnames]{xcolor}
\usepackage{enumitem}
\usepackage{listings}
\usepackage{tcolorbox}
\usepackage{titlesec}
\usepackage{fancyhdr}
\usepackage{array}
\usepackage{longtable}
\usepackage{booktabs}
\usepackage{amsmath}
\usepackage{amssymb}
\usepackage{hyperref}
\usepackage{cleveref}

\tcbuselibrary{skins,breakable}

\definecolor{specBlue}{RGB}{30,64,132}
\definecolor{specGray}{RGB}{90,90,90}
\definecolor{specAccent}{RGB}{198,40,40}
\definecolor{specGreen}{RGB}{30,110,60}

\hypersetup{
  colorlinks=true,
  linkcolor=specBlue,
  urlcolor=specBlue,
  citecolor=specBlue,
  pdfborder={0 0 0}
}

\titleformat{\section}{\Large\bfseries\color{specBlue}}{\thesection}{0.75em}{}
\titleformat{\subsection}{\large\bfseries\color{specBlue!80}}{\thesubsection}{0.6em}{}

\makeatletter
\newcommand{\specID}[1]{\def\@specID{#1}}
\newcommand{\specTitle}[1]{\def\@specTitle{#1}}
\newcommand{\specStatus}[1]{\def\@specStatus{#1}}
\newcommand{\specVersion}[1]{\def\@specVersion{#1}}
\newcommand{\specOwner}[1]{\def\@specOwner{#1}}
\newcommand{\specTraces}[1]{\def\@specTraces{#1}}

\specID{SPEC-XXX}
\specTitle{Untitled Spec}
\specStatus{DRAFT}
\specVersion{0.1}
\specOwner{Jeremie Nunez}
\specTraces{}

\newcommand{\makespecheader}{%
  \begin{center}
    {\LARGE\bfseries\color{specBlue}\@specTitle}\\[0.4em]
    {\small\color{specGray}\texttt{\@specID} \quad v\@specVersion \quad \textsc{\@specStatus} \quad \@specOwner}
  \end{center}
  \vspace{0.4em}
  \hrule
  \vspace{0.8em}
}

\pagestyle{fancy}
\fancyhf{}
\fancyhead[L]{\small\color{specGray}\@specID}
\fancyhead[R]{\small\color{specGray}\@specTitle}
\fancyfoot[C]{\small\color{specGray}\thepage}
\renewcommand{\headrulewidth}{0pt}
\makeatother

\newtcolorbox{invariant}{
  colback=specBlue!5, colframe=specBlue!75,
  title=Invariant, fonttitle=\bfseries, breakable,
  boxrule=0.5pt, arc=2pt
}

\newtcolorbox{scenario}[1]{
  colback=white, colframe=specBlue!50,
  title={Scenario: #1}, fonttitle=\bfseries, breakable,
  boxrule=0.5pt, arc=2pt
}

\newtcolorbox{ac}[1]{
  colback=specAccent!5, colframe=specAccent!60,
  title={Acceptance Criterion #1}, fonttitle=\bfseries, breakable,
  boxrule=0.5pt, arc=2pt
}

\newtcolorbox{nongoal}{
  colback=specGray!8, colframe=specGray!60,
  title=Non-Goal, fonttitle=\bfseries, breakable,
  boxrule=0.5pt, arc=2pt
}

\lstdefinestyle{codestyle}{
  basicstyle=\ttfamily\small,
  breaklines=true,
  showstringspaces=false,
  frame=single,
  framesep=4pt,
  rulecolor=\color{specBlue!30},
  backgroundcolor=\color{specBlue!2},
  xleftmargin=0.5em,
  xrightmargin=0.5em,
  keywordstyle=\color{specBlue}\bfseries,
  commentstyle=\color{specGray}\itshape,
  stringstyle=\color{specGreen}
}
\lstset{
  inputencoding=utf8,
  extendedchars=true,
  literate=
    {—}{{--}}1
    {–}{{-}}1
    {…}{{...}}1
    {’}{{'}}1
    {‘}{{'}}1
    {“}{{``}}1
    {”}{{''}}1
    {é}{{\'e}}1
    {è}{{\`e}}1
    {ê}{{\^e}}1
    {à}{{\`a}}1
    {â}{{\^a}}1
    {ç}{{\c{c}}}1
    {ô}{{\^o}}1
    {→}{{$\to$}}1
    {←}{{$\leftarrow$}}1
    {×}{{$\times$}}1
    {≥}{{$\geq$}}1
    {≤}{{$\leq$}}1
    {≠}{{$\neq$}}1
    {∈}{{$\in$}}1
    {⌈}{{$\lceil$}}1
    {⌉}{{$\rceil$}}1
    {∞}{{$\infty$}}1
    {α}{{$\alpha$}}1
    {β}{{$\beta$}}1
}
\lstdefinelanguage{TypeScript}{
  keywords={const,let,var,function,return,if,else,for,while,class,interface,type,extends,implements,import,export,from,as,async,await,new,this,true,false,null,undefined,void,any,unknown,never,string,number,boolean,readonly,private,public,protected,enum},
  sensitive=true,
  comment=[l]{//},
  morecomment=[s]{/*}{*/},
  morestring=[b]",
  morestring=[b]',
  morestring=[b]`
}
\lstset{style=codestyle}

\newcommand{\Given}{\textbf{\color{specBlue}Given} }
\newcommand{\When}{\textbf{\color{specBlue}When} }
\newcommand{\Then}{\textbf{\color{specBlue}Then} }
\providecommand{\And}{}
\renewcommand{\And}{\textbf{\color{specBlue}And} }
 (adjust depth).

\documentclass[11pt,a4paper]{article}

\usepackage[utf8]{inputenc}
\usepackage[T1]{fontenc}
\usepackage{lmodern}
\usepackage[a4paper,margin=2.3cm]{geometry}
\usepackage{parskip}
\usepackage{microtype}
\usepackage[dvipsnames]{xcolor}
\usepackage{enumitem}
\usepackage{listings}
\usepackage{tcolorbox}
\usepackage{titlesec}
\usepackage{fancyhdr}
\usepackage{array}
\usepackage{longtable}
\usepackage{booktabs}
\usepackage{amsmath}
\usepackage{amssymb}
\usepackage{hyperref}
\usepackage{cleveref}

\tcbuselibrary{skins,breakable}

\definecolor{specBlue}{RGB}{30,64,132}
\definecolor{specGray}{RGB}{90,90,90}
\definecolor{specAccent}{RGB}{198,40,40}
\definecolor{specGreen}{RGB}{30,110,60}

\hypersetup{
  colorlinks=true,
  linkcolor=specBlue,
  urlcolor=specBlue,
  citecolor=specBlue,
  pdfborder={0 0 0}
}

\titleformat{\section}{\Large\bfseries\color{specBlue}}{\thesection}{0.75em}{}
\titleformat{\subsection}{\large\bfseries\color{specBlue!80}}{\thesubsection}{0.6em}{}

\makeatletter
\newcommand{\specID}[1]{\def\@specID{#1}}
\newcommand{\specTitle}[1]{\def\@specTitle{#1}}
\newcommand{\specStatus}[1]{\def\@specStatus{#1}}
\newcommand{\specVersion}[1]{\def\@specVersion{#1}}
\newcommand{\specOwner}[1]{\def\@specOwner{#1}}
\newcommand{\specTraces}[1]{\def\@specTraces{#1}}

\specID{SPEC-XXX}
\specTitle{Untitled Spec}
\specStatus{DRAFT}
\specVersion{0.1}
\specOwner{Jeremie Nunez}
\specTraces{}

\newcommand{\makespecheader}{%
  \begin{center}
    {\LARGE\bfseries\color{specBlue}\@specTitle}\\[0.4em]
    {\small\color{specGray}\texttt{\@specID} \quad v\@specVersion \quad \textsc{\@specStatus} \quad \@specOwner}
  \end{center}
  \vspace{0.4em}
  \hrule
  \vspace{0.8em}
}

\pagestyle{fancy}
\fancyhf{}
\fancyhead[L]{\small\color{specGray}\@specID}
\fancyhead[R]{\small\color{specGray}\@specTitle}
\fancyfoot[C]{\small\color{specGray}\thepage}
\renewcommand{\headrulewidth}{0pt}
\makeatother

\newtcolorbox{invariant}{
  colback=specBlue!5, colframe=specBlue!75,
  title=Invariant, fonttitle=\bfseries, breakable,
  boxrule=0.5pt, arc=2pt
}

\newtcolorbox{scenario}[1]{
  colback=white, colframe=specBlue!50,
  title={Scenario: #1}, fonttitle=\bfseries, breakable,
  boxrule=0.5pt, arc=2pt
}

\newtcolorbox{ac}[1]{
  colback=specAccent!5, colframe=specAccent!60,
  title={Acceptance Criterion #1}, fonttitle=\bfseries, breakable,
  boxrule=0.5pt, arc=2pt
}

\newtcolorbox{nongoal}{
  colback=specGray!8, colframe=specGray!60,
  title=Non-Goal, fonttitle=\bfseries, breakable,
  boxrule=0.5pt, arc=2pt
}

\lstdefinestyle{codestyle}{
  basicstyle=\ttfamily\small,
  breaklines=true,
  showstringspaces=false,
  frame=single,
  framesep=4pt,
  rulecolor=\color{specBlue!30},
  backgroundcolor=\color{specBlue!2},
  xleftmargin=0.5em,
  xrightmargin=0.5em,
  keywordstyle=\color{specBlue}\bfseries,
  commentstyle=\color{specGray}\itshape,
  stringstyle=\color{specGreen}
}
\lstset{
  inputencoding=utf8,
  extendedchars=true,
  literate=
    {—}{{--}}1
    {–}{{-}}1
    {…}{{...}}1
    {’}{{'}}1
    {‘}{{'}}1
    {“}{{``}}1
    {”}{{''}}1
    {é}{{\'e}}1
    {è}{{\`e}}1
    {ê}{{\^e}}1
    {à}{{\`a}}1
    {â}{{\^a}}1
    {ç}{{\c{c}}}1
    {ô}{{\^o}}1
    {→}{{$\to$}}1
    {←}{{$\leftarrow$}}1
    {×}{{$\times$}}1
    {≥}{{$\geq$}}1
    {≤}{{$\leq$}}1
    {≠}{{$\neq$}}1
    {∈}{{$\in$}}1
    {⌈}{{$\lceil$}}1
    {⌉}{{$\rceil$}}1
    {∞}{{$\infty$}}1
    {α}{{$\alpha$}}1
    {β}{{$\beta$}}1
}
\lstdefinelanguage{TypeScript}{
  keywords={const,let,var,function,return,if,else,for,while,class,interface,type,extends,implements,import,export,from,as,async,await,new,this,true,false,null,undefined,void,any,unknown,never,string,number,boolean,readonly,private,public,protected,enum},
  sensitive=true,
  comment=[l]{//},
  morecomment=[s]{/*}{*/},
  morestring=[b]",
  morestring=[b]',
  morestring=[b]`
}
\lstset{style=codestyle}

\newcommand{\Given}{\textbf{\color{specBlue}Given} }
\newcommand{\When}{\textbf{\color{specBlue}When} }
\newcommand{\Then}{\textbf{\color{specBlue}Then} }
\providecommand{\And}{}
\renewcommand{\And}{\textbf{\color{specBlue}And} }
 (adjust depth).

\documentclass[11pt,a4paper]{article}

\usepackage[utf8]{inputenc}
\usepackage[T1]{fontenc}
\usepackage{lmodern}
\usepackage[a4paper,margin=2.3cm]{geometry}
\usepackage{parskip}
\usepackage{microtype}
\usepackage[dvipsnames]{xcolor}
\usepackage{enumitem}
\usepackage{listings}
\usepackage{tcolorbox}
\usepackage{titlesec}
\usepackage{fancyhdr}
\usepackage{array}
\usepackage{longtable}
\usepackage{booktabs}
\usepackage{amsmath}
\usepackage{amssymb}
\usepackage{hyperref}
\usepackage{cleveref}

\tcbuselibrary{skins,breakable}

\definecolor{specBlue}{RGB}{30,64,132}
\definecolor{specGray}{RGB}{90,90,90}
\definecolor{specAccent}{RGB}{198,40,40}
\definecolor{specGreen}{RGB}{30,110,60}

\hypersetup{
  colorlinks=true,
  linkcolor=specBlue,
  urlcolor=specBlue,
  citecolor=specBlue,
  pdfborder={0 0 0}
}

\titleformat{\section}{\Large\bfseries\color{specBlue}}{\thesection}{0.75em}{}
\titleformat{\subsection}{\large\bfseries\color{specBlue!80}}{\thesubsection}{0.6em}{}

\makeatletter
\newcommand{\specID}[1]{\def\@specID{#1}}
\newcommand{\specTitle}[1]{\def\@specTitle{#1}}
\newcommand{\specStatus}[1]{\def\@specStatus{#1}}
\newcommand{\specVersion}[1]{\def\@specVersion{#1}}
\newcommand{\specOwner}[1]{\def\@specOwner{#1}}
\newcommand{\specTraces}[1]{\def\@specTraces{#1}}

\specID{SPEC-XXX}
\specTitle{Untitled Spec}
\specStatus{DRAFT}
\specVersion{0.1}
\specOwner{Jeremie Nunez}
\specTraces{}

\newcommand{\makespecheader}{%
  \begin{center}
    {\LARGE\bfseries\color{specBlue}\@specTitle}\\[0.4em]
    {\small\color{specGray}\texttt{\@specID} \quad v\@specVersion \quad \textsc{\@specStatus} \quad \@specOwner}
  \end{center}
  \vspace{0.4em}
  \hrule
  \vspace{0.8em}
}

\pagestyle{fancy}
\fancyhf{}
\fancyhead[L]{\small\color{specGray}\@specID}
\fancyhead[R]{\small\color{specGray}\@specTitle}
\fancyfoot[C]{\small\color{specGray}\thepage}
\renewcommand{\headrulewidth}{0pt}
\makeatother

\newtcolorbox{invariant}{
  colback=specBlue!5, colframe=specBlue!75,
  title=Invariant, fonttitle=\bfseries, breakable,
  boxrule=0.5pt, arc=2pt
}

\newtcolorbox{scenario}[1]{
  colback=white, colframe=specBlue!50,
  title={Scenario: #1}, fonttitle=\bfseries, breakable,
  boxrule=0.5pt, arc=2pt
}

\newtcolorbox{ac}[1]{
  colback=specAccent!5, colframe=specAccent!60,
  title={Acceptance Criterion #1}, fonttitle=\bfseries, breakable,
  boxrule=0.5pt, arc=2pt
}

\newtcolorbox{nongoal}{
  colback=specGray!8, colframe=specGray!60,
  title=Non-Goal, fonttitle=\bfseries, breakable,
  boxrule=0.5pt, arc=2pt
}

\lstdefinestyle{codestyle}{
  basicstyle=\ttfamily\small,
  breaklines=true,
  showstringspaces=false,
  frame=single,
  framesep=4pt,
  rulecolor=\color{specBlue!30},
  backgroundcolor=\color{specBlue!2},
  xleftmargin=0.5em,
  xrightmargin=0.5em,
  keywordstyle=\color{specBlue}\bfseries,
  commentstyle=\color{specGray}\itshape,
  stringstyle=\color{specGreen}
}
\lstset{
  inputencoding=utf8,
  extendedchars=true,
  literate=
    {—}{{--}}1
    {–}{{-}}1
    {…}{{...}}1
    {’}{{'}}1
    {‘}{{'}}1
    {“}{{``}}1
    {”}{{''}}1
    {é}{{\'e}}1
    {è}{{\`e}}1
    {ê}{{\^e}}1
    {à}{{\`a}}1
    {â}{{\^a}}1
    {ç}{{\c{c}}}1
    {ô}{{\^o}}1
    {→}{{$\to$}}1
    {←}{{$\leftarrow$}}1
    {×}{{$\times$}}1
    {≥}{{$\geq$}}1
    {≤}{{$\leq$}}1
    {≠}{{$\neq$}}1
    {∈}{{$\in$}}1
    {⌈}{{$\lceil$}}1
    {⌉}{{$\rceil$}}1
    {∞}{{$\infty$}}1
    {α}{{$\alpha$}}1
    {β}{{$\beta$}}1
}
\lstdefinelanguage{TypeScript}{
  keywords={const,let,var,function,return,if,else,for,while,class,interface,type,extends,implements,import,export,from,as,async,await,new,this,true,false,null,undefined,void,any,unknown,never,string,number,boolean,readonly,private,public,protected,enum},
  sensitive=true,
  comment=[l]{//},
  morecomment=[s]{/*}{*/},
  morestring=[b]",
  morestring=[b]',
  morestring=[b]`
}
\lstset{style=codestyle}

\newcommand{\Given}{\textbf{\color{specBlue}Given} }
\newcommand{\When}{\textbf{\color{specBlue}When} }
\newcommand{\Then}{\textbf{\color{specBlue}Then} }
\providecommand{\And}{}
\renewcommand{\And}{\textbf{\color{specBlue}And} }


\specID{SPEC-SW-006}
\specTitle{Multi-Agent Simulation Engine --- operator behaviour \& conjunction negotiation}
\specStatus{DRAFT}
\specVersion{0.1}
\specOwner{Jeremie Nunez}

\begin{document}
\makespecheader

\section{Context}
Thalamus currently answers \emph{what is} and \emph{what just happened} in the SSA knowledge graph. It cannot answer \emph{what would happen if}. This spec introduces a multi-agent simulation engine, hosted in \texttt{packages/sweep}, that instantiates persistent agents from KG state, runs them through a turn-based loop, and produces a narrative report plus a deep-interaction chat surface.

Inspiration: MiroFish's five-stage pipeline (\emph{graph build $\to$ environment setup $\to$ simulation $\to$ report $\to$ deep interaction}). This spec adapts that pattern to SSA, reusing the existing cortex registry, LLM transport, pgvector store, and BullMQ queues rather than running a Python sidecar.

\subsection*{Core idea --- the swarm, not the sim}
The primary unit of value is NOT a single high-fidelity simulation. It is a \textbf{swarm of cheap small-model ``fish''} that collectively cover the space of plausible outcomes. Each fish is a short, nano-model-driven mini-sim with a perturbed seed (different god-event, different initial assumption, different persona tuning). The aggregator clusters outcomes and produces a distribution --- ``in 62\% of futures, operator A maneuvers; in 28\%, B maneuvers; in 10\%, negotiation fails and escalates''.

A single run is just a swarm of size 1. Determinism guarantees mean fixture mode can replay the full swarm byte-identically.

Two use cases are in-scope for v0:
\begin{itemize}[leftmargin=1.2em]
  \item \textbf{UC1 --- Operator behaviour swarm.} $N$ agents (one per operator), run for a short compressed horizon, across $K$ fish ($K \approx 50$--$100$) each with a different god-event perturbation (regulation variants, ASAT event on different satellites, launch surges at different scales, debris cascade severities). Output: a coverage report over operator decisions --- ``under these perturbations, what fraction of futures trigger a Starlink slot-share reshuffle?''.
  \item \textbf{UC3 --- Conjunction negotiation swarm.} A 2-agent swarm triggered by a conjunction finding. $K$ fish ($K \approx 30$--$50$) each perturb the seed (delta-v budget, risk tolerance, insurance-loss assumption). Output: a distribution over resolutions (who maneuvers, split ratio) plus a confidence band. Feeds the modal resolution back into \texttt{finding-routing} as a \texttt{sweep\_suggestion}, annotated with the distribution for reviewer context.
\end{itemize}

A fish runs on the nano model (\texttt{gpt-5.4-nano} via \texttt{callNano}) by default; swarm cost is $\sim\$0.01$/fish, so $K = 100$ fish costs about \$1.

\section{Goals}
\begin{itemize}[leftmargin=1.2em]
  \item \textbf{Swarm-first}: the orchestrator's primary API is \texttt{launchSwarm}; a single sim is $K=1$. Each fish is a cheap nano-driven mini-sim with a perturbed seed.
  \item \textbf{Coverage over fidelity}: prefer 100 short cheap fish to 1 long expensive sim. The aggregator produces a distribution of outcomes, not a verdict.
  \item Model long-lived \textbf{agents within each fish} (not one-shot cortex calls) with stable identity, persona, and memory across turns.
  \item Drive \textbf{fish execution} over BullMQ with parallel fan-out (concurrency 8--16) so a 100-fish swarm finishes in $<$ 2 minutes wall-clock on fixtures.
  \item Persist \textbf{per-agent memory} in pgvector scoped by \texttt{(sim\_run\_id, agent\_id)} so fish cannot bleed into each other.
  \item Reuse existing LLM pathways: \texttt{callNano} (nano default, cheap) for fish bodies; \texttt{createLlmTransportWithMode} (Kimi K2) for the aggregator / reporter.
  \item Reuse the \textbf{cortex registry} as the per-turn reasoning primitive --- no new LLM abstraction.
  \item Produce a \textbf{swarm coverage report} via a new \texttt{sim\_swarm\_reporter} cortex; a single-fish report via \texttt{sim\_reporter}.
  \item Wire \textbf{UC3} into \texttt{finding-routing} so conjunction findings can automatically spawn a negotiation \emph{swarm} whose modal resolution feeds back as a sweep suggestion, annotated with the full distribution for reviewer context.
\end{itemize}

\section{Non-Goals}
\begin{nongoal}
This spec does NOT re-implement orbital propagation, does NOT compute real delta-v budgets, and does NOT produce legally-actionable maneuver plans. All physics is heuristic: mass-based cost bands (already available via \texttt{queryReplacementCost}) and fleet-level proxies. Realism is narrative, not numerical.
\end{nongoal}
\begin{nongoal}
This spec does NOT introduce streaming LLM calls inside the sim loop. Each turn is a single \texttt{await} against \texttt{LlmChatTransport}; streaming remains confined to the existing chat surface.
\end{nongoal}
\begin{nongoal}
This spec does NOT require MiroFish's Python runtime, Zep Cloud, or OASIS. The only runtime additions are (a) two new Postgres tables and (b) two new BullMQ workers.
\end{nongoal}

\section{Invariants}
\begin{invariant}
Agents are append-only. An agent's \texttt{agent.id} persists across sim runs; \texttt{agent\_memory} rows are never mutated, only appended and soft-retired via \texttt{retired\_at}.
\end{invariant}
\begin{invariant}
A \textbf{turn} is the atomic unit of sim progress. Partial turns are rolled back: if any agent's turn call fails, the entire turn is discarded and BullMQ retries with exponential backoff (max 3 attempts).
\end{invariant}
\begin{invariant}
God-view injections (policy, event) are themselves \texttt{sim\_turn} rows with \texttt{actor\_kind = 'god'} --- they live in the same timeline as agent turns so the narrative report can interleave them.
\end{invariant}
\begin{invariant}
No cross-sim memory bleed. An agent bound to operator \texttt{X} in \texttt{sim\_run} \texttt{A} cannot read memory rows written in \texttt{sim\_run} \texttt{B}. Memory is scoped by \texttt{(agent\_id, sim\_run\_id)}.
\end{invariant}
\begin{invariant}
UC3 negotiation swarms produce at most one \texttt{sweep\_suggestion} per swarm (not per fish), tagged with \texttt{sim\_swarm\_id}. The suggestion payload carries the modal resolution plus the full outcome distribution as metadata. Individual fish never emit suggestions directly.
\end{invariant}
\begin{invariant}
A swarm is deterministic under fixed \texttt{(baseSeed, perturbationSet, size, llmMode="fixtures")}. Re-running the same swarm config yields byte-identical fish outcomes and byte-identical aggregator output.
\end{invariant}
\begin{invariant}
Fish within a swarm are \emph{independent} --- no fish reads another fish's memory or timeline. Cross-fish aggregation happens only in the aggregator, which reads \texttt{sim\_turn} rows grouped by \texttt{sim\_run\_id}. This guarantees embarrassingly-parallel fan-out.
\end{invariant}
\begin{invariant}
Swarm abort is fail-soft: if $f$ of $K$ fish fail (LLM error, timeout), the aggregator proceeds when $K - f \geq \lceil 0.8 K \rceil$ (80\% quorum). Below quorum, the swarm transitions to \texttt{failed} and no suggestion is emitted. Per-fish timeout is 60s.
\end{invariant}
\begin{invariant}
Thalamus never imports sweep. All new code lives in \texttt{packages/sweep/src/sim/}; the only Thalamus surfaces consumed are \texttt{CortexExecutor.runSkillFreeform}, \texttt{ThalamusDAGExecutor}, \texttt{createLlmTransportWithMode}, and the existing embedding helper. The two new cortex skills (\texttt{sim\_operator\_agent}, \texttt{sim\_reporter}) live in Thalamus and are consumed from sweep through the standard cortex registry --- no special coupling.
\end{invariant}
\begin{invariant}
\texttt{sim\_turn} is the canonical, mode-agnostic timeline. Both drivers (DAG and Sequential) write every agent action and every god event there. The reporter and the chat surface query \texttt{sim\_turn} exclusively --- never the DAG driver's transient \texttt{research\_finding} rows --- so the two modes are observationally indistinguishable.
\end{invariant}
\begin{invariant}
Promotion of a \texttt{sim\_turn} row to a \texttt{research\_finding} is action-kind gated: only \texttt{maneuver}, \texttt{launch}, \texttt{retire} are promoted. Negotiation micro-actions (\texttt{propose\_split}, \texttt{accept}, \texttt{reject}, \texttt{hold}, \texttt{lobby}) MUST NOT create \texttt{research\_finding} rows. Rationale: keep the KG signal-rich; negotiation noise stays scoped to the sim run.
\end{invariant}

\section{Architecture}

\subsection{Swarm layer (top)}
The swarm is the public API. It fans out $K$ fish over BullMQ, each fish is one \texttt{sim\_run}, each fish picks a driver (DAG or Sequential) by its kind. The aggregator reads \texttt{sim\_turn} across all fish in the swarm and emits a coverage report.

\begin{lstlisting}[language=TypeScript]
// swarm.service.ts
launchSwarm({
  kind: "uc1_operator_behavior" | "uc3_conjunction",
  baseSeed: SeedRefs,                    // fleet, finding, horizon
  perturbations: PerturbationSpec[],     // K entries
  llmMode: "cloud" | "fixtures" | "record",
}): Promise<{ swarmId: number }>
\end{lstlisting}

Each \texttt{PerturbationSpec} is applied to the base seed to produce a fish-specific seed. Examples:
\begin{itemize}[leftmargin=1.2em]
  \item \texttt{\{ godEvent: \{ kind: "asat\_event", targetSatelliteId: 42 \} \}}
  \item \texttt{\{ constraintOverride: \{ agentId: 7, maxDeltaVBudgetPerSat: 50 \} \}} (tighter budget for one agent)
  \item \texttt{\{ personaTweak: \{ agentId: 7, riskProfile: "aggressive" \} \}}
  \item \texttt{\{ launchSurge: \{ regimeId: 3, extraSatellites: 100 \} \}}
  \item \texttt{\{ noop: true \}} (baseline control fish, always present as fish[0])
\end{itemize}

\subsection{Two drivers, one fish}
The sim engine ships with \textbf{two turn drivers} selected by \texttt{sim\_run.kind}:

\begin{itemize}[leftmargin=1.2em]
  \item \textbf{DAG driver} (UC1) --- one turn = one \texttt{ThalamusDAGExecutor.execute()} call. Agents are parallel cortex nodes (\texttt{sim\_operator\_agent}) in the same level; a reconciler node settles observations. Each turn is a \texttt{research\_cycle}, giving a full audit trail for free (findings, edges, embeddings). Optimal when $N$ agents decide simultaneously.
  \item \textbf{Sequential driver} (UC3) --- one turn = one direct \texttt{CortexExecutor.runSkillFreeform()} call against the same \texttt{sim\_operator\_agent} skill. Agents alternate (A $\to$ B $\to$ A $\ldots$), intermediate state (current offer) lives only in \texttt{sim\_turn}, no \texttt{research\_finding} noise. Optimal for tight negotiation back-and-forth.
\end{itemize}

Both drivers write to the same \texttt{sim\_turn} table --- the canonical, mode-agnostic timeline consumed by the reporter and the chat surface. The cortex skill \texttt{sim\_operator\_agent} is shared: one prompt, two invocation paths.

\paragraph{Research-finding promotion.} Every agent action is written to \texttt{sim\_turn}. A subset (\texttt{action.kind $\in$ \{maneuver, launch, retire\}}) is additionally promoted to \texttt{research\_finding} for KG-level audit (embedding, edges to satellite / operator). Negotiation micro-actions (\texttt{propose\_split, accept, reject, hold, lobby}) remain in \texttt{sim\_turn} only, so the KG stays signal-rich.

\subsection{Package layout}
\begin{lstlisting}[language=TypeScript]
packages/sweep/src/sim/
  types.ts                      // SimSwarm, SimRun, Agent, Turn, Memory, GodEvent, Perturbation
  schema.ts                     // Zod schemas for turn outputs + perturbations
  agent-builder.ts              // KG -> Agent (persona synthesis)
  memory.service.ts             // write/read agent_memory (pgvector, scoped)
  swarm.service.ts              // fan-out K fish, track quorum
  perturbation.ts               // generators + applicator (seed + spec -> fish seed)
  aggregator.service.ts         // cluster fish outcomes, compute distribution
  turn-runner-dag.ts            // UC1 fish driver: build DAG, runCycle, persist
  turn-runner-sequential.ts     // UC3 fish driver: direct runSkillFreeform
  god-channel.service.ts        // inject events (per-fish, pre-seeded via perturbation)
  sim-orchestrator.service.ts   // internal: schedules turns within a fish
  reporter.service.ts           // renders single-fish narrative
  swarm-reporter.service.ts     // renders swarm coverage report
  promote.ts                    // modal outcome -> research_finding + sweep_suggestion

packages/sweep/src/jobs/workers/
  sim-turn.worker.ts            // one agent-turn, picks driver by sim_run.kind
  swarm-fish.worker.ts          // one fish lifecycle (start -> drain turns -> close)
  swarm-aggregate.worker.ts     // fires when quorum reached -> aggregator + reporter

packages/sweep/src/routes/
  swarm.routes.ts               // POST /swarm/uc1, POST /swarm/uc3, GET /swarm/:id

packages/thalamus/src/cortices/skills/
  sim_operator_agent.md         // NEW: shared per-turn agent prompt (nano-friendly)
  sim_reconciler.md             // NEW: no-LLM, post-turn settler for DAG driver
  sim_reporter.md               // NEW: single-fish narrative
  sim_swarm_reporter.md         // NEW: swarm-level coverage report

packages/thalamus/src/cortices/
  sql-helpers.sim-context.ts    // querySimAgentContext(simRunId, agentId, turn)
  sql-helpers.sim-timeline.ts   // querySimTimeline(simRunId) for single-fish reporter
  sql-helpers.sim-swarm.ts      // querySwarmOutcomes(swarmId) for aggregator + reporter
\end{lstlisting}

\subsection{Data model}
Five tables in \texttt{packages/db-schema/src/schema/sim.ts}:

\begin{lstlisting}[language=TypeScript]
export const simSwarm = pgTable("sim_swarm", {
  id: bigserial("id", { mode: "number" }).primaryKey(),
  kind: text("kind").$type<"uc1_operator_behavior" | "uc3_conjunction">().notNull(),
  title: text("title").notNull(),
  baseSeed: jsonb("base_seed").$type<SeedRefs>().notNull(),
  perturbations: jsonb("perturbations").$type<PerturbationSpec[]>().notNull(),
  size: integer("size").notNull(),             // = perturbations.length
  config: jsonb("config").$type<SwarmConfig>().notNull(),
  // { llmMode, quorumPct: 0.8, perFishTimeoutMs: 60000, fishConcurrency: 8 }
  status: text("status").$type<"pending"|"running"|"done"|"failed">().notNull(),
  outcomeReportFindingId: bigint("outcome_report_finding_id", { mode: "number" }),
  suggestionId: bigint("suggestion_id", { mode: "number" }), // UC3 modal outcome
  startedAt: timestamp("started_at", { withTimezone: true }).defaultNow(),
  completedAt: timestamp("completed_at", { withTimezone: true }),
  createdBy: bigint("created_by", { mode: "number" }),
});

export const simRun = pgTable("sim_run", {
  id: bigserial("id", { mode: "number" }).primaryKey(),
  swarmId: bigint("swarm_id", { mode: "number" })
    .references(() => simSwarm.id, { onDelete: "cascade" }).notNull(),
  fishIndex: integer("fish_index").notNull(),  // 0..swarm.size - 1
  kind: text("kind").$type<"uc1_operator_behavior" | "uc3_conjunction">().notNull(),
  title: text("title").notNull(),
  seedRefs: jsonb("seed_refs").$type<SeedRefs>().notNull(),
  // e.g. { operatorIds: [...], conjunctionFindingId?: number, horizonDays: 30 }
  config: jsonb("config").$type<SimConfig>().notNull(),
  // { turnsPerDay: 1, maxTurns: 90, llmMode: "fixtures"|"cloud", seed: 42 }
  status: text("status").$type<"pending"|"running"|"paused"|"done"|"failed">().notNull(),
  startedAt: timestamp("started_at", { withTimezone: true }).defaultNow(),
  completedAt: timestamp("completed_at", { withTimezone: true }),
  reportFindingId: bigint("report_finding_id", { mode: "number" }),
  createdBy: bigint("created_by", { mode: "number" }),
});

export const simAgent = pgTable("sim_agent", {
  id: bigserial("id", { mode: "number" }).primaryKey(),
  simRunId: bigint("sim_run_id", { mode: "number" })
    .references(() => simRun.id, { onDelete: "cascade" }).notNull(),
  operatorId: bigint("operator_id", { mode: "number" })
    .references(() => operator.id),
  persona: text("persona").notNull(), // system-prompt fragment
  goals: jsonb("goals").$type<string[]>().notNull(),
  constraints: jsonb("constraints").$type<Record<string, unknown>>().notNull(),
  createdAt: timestamp("created_at", { withTimezone: true }).defaultNow(),
});

export const simTurn = pgTable("sim_turn", {
  id: bigserial("id", { mode: "number" }).primaryKey(),
  simRunId: bigint("sim_run_id", { mode: "number" })
    .references(() => simRun.id, { onDelete: "cascade" }).notNull(),
  turnIndex: integer("turn_index").notNull(),
  actorKind: text("actor_kind").$type<"agent"|"god"|"system">().notNull(),
  agentId: bigint("agent_id", { mode: "number" })
    .references(() => simAgent.id),
  action: jsonb("action").$type<TurnAction>().notNull(),
  rationale: text("rationale").notNull(),
  observableSummary: text("observable_summary").notNull(),
  // what other agents see next turn; may omit private rationale
  llmCostUsd: numeric("llm_cost_usd", { precision: 10, scale: 6 }),
  createdAt: timestamp("created_at", { withTimezone: true }).defaultNow(),
}, (t) => ({
  uniqTurn: uniqueIndex("sim_turn_uniq").on(t.simRunId, t.turnIndex, t.agentId),
}));

export const simAgentMemory = pgTable("sim_agent_memory", {
  id: bigserial("id", { mode: "number" }).primaryKey(),
  simRunId: bigint("sim_run_id", { mode: "number" })
    .references(() => simRun.id, { onDelete: "cascade" }).notNull(),
  agentId: bigint("agent_id", { mode: "number" })
    .references(() => simAgent.id).notNull(),
  turnIndex: integer("turn_index").notNull(),
  kind: text("kind").$type<"self_action"|"observation"|"belief">().notNull(),
  content: text("content").notNull(),
  embedding: vector("embedding", { dimensions: EMBEDDING_DIMENSIONS }),
  retiredAt: timestamp("retired_at", { withTimezone: true }),
});
\end{lstlisting}

HNSW index on \texttt{sim\_agent\_memory.embedding} follows the existing pattern (\texttt{migrations/0001\_hnsw\_index.sql}).

\subsection{Turn loop --- DAG driver (UC1)}

\begin{lstlisting}[language=TypeScript]
// turn-runner-dag.ts --- one turn = one DAG cycle, N agents in parallel
async function runTurnDag(simRunId: number, turnIndex: number) {
  const agents = await loadAgents(simRunId);

  const dag: DAGPlan = {
    intent: `sim ${simRunId} turn ${turnIndex}`,
    complexity: "moderate",
    nodes: [
      ...agents.map((a) => ({
        cortex: "sim_operator_agent",
        params: { simRunId, agentId: a.id, turnIndex },
        dependsOn: [],
      })),
      {
        cortex: "sim_reconciler",   // settles observations, writes sim_turn rows
        params: { simRunId, turnIndex },
        dependsOn: agents.map((_, i) => `sim_operator_agent#${i}`),
      },
    ],
  };

  // Existing Thalamus plumbing: persists research_cycle, findings, edges
  await thalamusService.runCycle({ dag, mode: llmMode });

  // Reconciler already wrote sim_turn + memory; we just advance.
  await scheduleNext(simRunId, turnIndex + 1);
}
\end{lstlisting}

\subsection{Turn loop --- Sequential driver (UC3)}

\begin{lstlisting}[language=TypeScript]
// turn-runner-sequential.ts --- one turn = one agent speaks, strict A->B alternation
async function runTurnSequential(simRunId: number, turnIndex: number) {
  const agent = pickSpeaker(simRunId, turnIndex); // alternates on turnIndex parity
  const ctx = await buildAgentContext(db, agent.id, turnIndex);

  const response = await cortexExecutor.runSkillFreeform({
    skillName: "sim_operator_agent",
    systemPrompt: agent.persona,
    userPrompt: renderTurnPrompt(ctx),
    schema: turnResponseSchema,
    mode: llmMode,
  });

  await db.transaction(async (tx) => {
    const turn = await insertSimTurn(tx, { simRunId, turnIndex, agent, response });
    await memory.write(tx, { self_action: turn });
    await memory.write(tx, { observation_for_others: turn.observableSummary });
    if (isKgPromotable(response.action)) {
      await promoteToFinding(tx, turn); // maneuver/launch/retire -> research_finding
    }
    if (isTerminal(response.action)) {
      await closeRun(tx, simRunId, turn); // accept / reject -> done
    }
  });

  await scheduleNext(simRunId, turnIndex + 1);
}
\end{lstlisting}

\subsection{God channel}

God events are injected via \texttt{POST /sim/:id/god} (admin-only). Payload is stored as a \texttt{sim\_turn} row with \texttt{actor\_kind = 'god'}. Agents see god events as part of their observable timeline --- no special branch in the agent prompt.

\subsection{UC3 wiring}

\texttt{finding-routing.ts} gains a new hook: if a finding has \texttt{cortex = 'conjunction\_analysis'} and \texttt{urgency >= 0.7}, enqueue a \texttt{sim-advance} job of kind \texttt{uc3\_conjunction} with the finding id as seed. The sim runs to completion (max 20 turns), its resolution becomes a \texttt{sweep\_suggestion} via \texttt{SweepRepository.enqueueSuggestion}, tagged with \texttt{sim\_run\_id}. Reviewer sees the suggestion in the normal inbox flow.

\section{Interface}

\begin{lstlisting}[language=TypeScript]
// sim-orchestrator.service.ts
export interface SimOrchestrator {
  startUc1(opts: {
    operatorIds: number[];
    horizonDays: number;
    turnsPerDay?: number;
    godEvents?: GodEventSeed[];
    mode?: "cloud" | "fixtures" | "record";
    createdBy?: number;
  }): Promise<{ simRunId: number }>;

  startUc3(opts: {
    conjunctionFindingId: number;
    maxTurns?: number;
    mode?: "cloud" | "fixtures" | "record";
  }): Promise<{ simRunId: number }>;

  pause(simRunId: number): Promise<void>;
  resume(simRunId: number): Promise<void>;
  inject(simRunId: number, event: GodEventSeed): Promise<void>;
  status(simRunId: number): Promise<SimStatus>;
}

// routes/sim.routes.ts
POST   /admin/sim/uc1         body: Uc1StartRequest
POST   /admin/sim/uc3         body: Uc3StartRequest
GET    /admin/sim/:id         -> SimStatus + last 20 turns
POST   /admin/sim/:id/god     body: GodEventSeed
POST   /admin/sim/:id/pause
POST   /admin/sim/:id/resume
GET    /admin/sim/:id/report  -> rendered narrative (from sim_reporter cortex)
\end{lstlisting}

\section{Scenarios}

\begin{scenario}{UC1 --- Free-running operator sim, 3 operators, 30-day horizon}
\Given three operators (SpaceX, OneWeb, CNSA) each with $\geq$ 50 satellites in KG \\
\And the orchestrator is booted with \texttt{mode = "fixtures"} for determinism \\
\When \texttt{startUc1(\{ operatorIds: [1,2,3], horizonDays: 30, turnsPerDay: 1 \})} is called \\
\Then a \texttt{sim\_run} row is created with \texttt{status = "running"} \\
\And three \texttt{sim\_agent} rows exist, one per operator, each with a persona synthesized from operator metadata + fleet stats \\
\And the BullMQ \texttt{sim-turn} queue has exactly 3 jobs queued (turn 0 for each agent) \\
\And after 90 turns (3 agents $\times$ 30 days), \texttt{status = "done"} and exactly 90 \texttt{sim\_turn} rows exist.
\end{scenario}

\begin{scenario}{UC1 --- God-view ASAT injection mid-run}
\Given a UC1 run is at turn 15 of 30, status \texttt{running} \\
\When admin POSTs \texttt{/admin/sim/:id/god} with \{\texttt{kind: "asat\_event"}, \texttt{targetSatelliteId: 42}\} \\
\Then one \texttt{sim\_turn} row is appended with \texttt{actor\_kind = "god"} \\
\And on turn 16, every agent's observable context includes the god event \\
\And at least one agent's turn-16 action references the ASAT event in its \texttt{rationale}.
\end{scenario}

\begin{scenario}{UC3 --- Conjunction auto-spawn from finding}
\Given a conjunction finding is persisted with cortex \texttt{conjunction\_analysis}, urgency $= 0.85$, and two satellites tied to operators A and B \\
\When \texttt{finding-routing} processes the finding \\
\Then a \texttt{sim\_run} of kind \texttt{uc3\_conjunction} is created within 2 seconds \\
\And two \texttt{sim\_agent} rows exist, one per operator \\
\And within 60 seconds the run completes with $\leq$ 20 turns \\
\And a \texttt{sweep\_suggestion} is enqueued with \texttt{sim\_run\_id} set \\
\And the suggestion's \texttt{actions} field includes a \texttt{maneuver} action targeting one of the two satellites.
\end{scenario}

\begin{scenario}{Fixture determinism}
\Given a UC3 run with \texttt{mode = "fixtures"}, \texttt{config.seed = 42}, fixed \texttt{conjunctionFindingId} \\
\When the run is executed twice in succession \\
\Then both runs produce byte-identical \texttt{sim\_turn.action} and \texttt{rationale} for every turn \\
\And both runs produce the same resolution \texttt{sweep\_suggestion} payload.
\end{scenario}

\begin{scenario}{Turn failure is atomic}
\Given a UC1 run where the LLM transport throws on turn 7 for agent 2 \\
\When the worker catches the error \\
\Then no \texttt{sim\_turn} row is written for turn 7 agent 2 \\
\And no \texttt{sim\_agent\_memory} row is written for that turn \\
\And BullMQ retries the job up to 3 times with exponential backoff \\
\And after 3 failures, the run transitions to \texttt{status = "failed"} and no further turns are scheduled.
\end{scenario}

\begin{scenario}{Memory scoping}
\Given operator A participated in sim run X (20 turns of memory) and sim run Y (5 turns) \\
\When agent-A in run Y queries its memory at turn 3 \\
\Then the vector search returns only rows where \texttt{sim\_run\_id = Y} \\
\And zero rows from run X are returned even if semantically similar.
\end{scenario}

\section{Acceptance criteria}

\begin{ac}{AC-1 Schema}
Migration applies cleanly: \texttt{make db-migrate} produces \texttt{sim\_run}, \texttt{sim\_agent}, \texttt{sim\_turn}, \texttt{sim\_agent\_memory} tables + HNSW index on \texttt{sim\_agent\_memory.embedding}. \texttt{pnpm typecheck} is 0 errors.
\end{ac}

\begin{ac}{AC-2 UC3 happy path}
\texttt{make sim-uc3} (new Makefile target) completes in $<$ 90s in fixtures mode, produces $\geq$ 4 turns, writes a \texttt{sweep\_suggestion} tagged with \texttt{sim\_run\_id}, and its \texttt{reportFindingId} resolves to a \texttt{research\_finding} row with non-empty summary.
\end{ac}

\begin{ac}{AC-3 UC1 happy path}
\texttt{make sim-uc1} with 3 operators and \texttt{horizonDays = 5} completes in $<$ 120s in fixtures mode, produces exactly 15 agent turns (ignoring god turns), and the final report mentions all 3 operators by name.
\end{ac}

\begin{ac}{AC-4 God injection}
\texttt{POST /admin/sim/:id/god} during a running UC1 appends a god turn within 500ms; the next agent turn's rationale references the god event string.
\end{ac}

\begin{ac}{AC-5 Determinism}
Two back-to-back fixture runs of UC3 with identical seed produce identical \texttt{action} JSON for every turn (byte-equal when sorted).
\end{ac}

\begin{ac}{AC-6 Isolation}
Thalamus package continues to typecheck and pass its existing tests with zero imports from \texttt{@interview/sweep}.
\end{ac}

\begin{ac}{AC-7 Chat deep-interaction}
The existing \texttt{satellite-sweep-chat} accepts a \texttt{simRunId} scope parameter; when set, its retrieval pulls from \texttt{sim\_turn} + \texttt{sim\_agent\_memory} for the given run instead of the global KG.
\end{ac}

\section{Open questions}
\begin{enumerate}[leftmargin=1.4em]
  \item Should \texttt{observableSummary} be generated by the same LLM call as \texttt{action + rationale}, or by a second cheap pass (\texttt{gpt-5.4-nano})? Two-pass is cleaner (prevents rationale leakage) but doubles cost.
  \item For UC1 at scale ($N \geq 10$ agents, $\geq 100$ days), should turns within a day run in parallel (all agents decide simultaneously, observations settle at end-of-day) or strictly sequential? Parallel is more realistic and faster; sequential is easier to debug.
  \item UC3: if the two agents fail to converge within \texttt{maxTurns}, should the sim emit a ``no-agreement'' suggestion (escalate to human reviewer) or a ``default: larger fleet maneuvers'' suggestion? Current proposal: escalate.
\end{enumerate}

\section{References}
\begin{itemize}[leftmargin=1.2em]
  \item MiroFish README --- \url{https://github.com/666ghj/MiroFish/blob/main/README.md}
  \item SPEC-SW-001 Nano Sweep (\texttt{nano-sweep.tex})
  \item SPEC-SW-002 Finding Routing (\texttt{finding-routing.tex}) --- UC3 integration point
  \item SPEC-SW-003 Resolution (\texttt{resolution.tex}) --- consumes UC3 output
  \item SPEC-TH-* Cortex Pattern --- \texttt{runSkillFreeform} is the per-turn primitive
\end{itemize}

\end{document}
